\documentclass[11pt]{article}

\usepackage[margin=1.2in]{geometry}
\usepackage{amsmath,amssymb}
\usepackage{microtype}
\usepackage{parskip}
\usepackage{listings}
\usepackage{xcolor}
\usepackage[hidelinks]{hyperref}
\usepackage{booktabs}
\usepackage{array}
\usepackage{enumitem}
\usepackage{mdframed}

\definecolor{codegray}{RGB}{245, 245, 245}
\definecolor{codecomment}{RGB}{100, 116, 139}
\definecolor{codestring}{RGB}{22, 101, 52}
\definecolor{codenumber}{RGB}{100, 116, 139}
\definecolor{bordercolor}{RGB}{203, 213, 225}
\definecolor{accentblue}{RGB}{30, 64, 175}
\definecolor{warningyellow}{RGB}{254, 243, 199}
\definecolor{warningborder}{RGB}{217, 119, 6}
\definecolor{notegreen}{RGB}{220, 252, 231}
\definecolor{noteborder}{RGB}{22, 101, 52}

\lstdefinelanguage{Rust}{
  keywords={pub, struct, enum, fn, let, mut, impl, use, self, Self,
            match, if, else, for, while, return, true, false,
            Vec, Option, Result, HashMap, HashSet, Some, None, Ok, Err,
            type, where, trait, as, in, ref, mod, super, Box, Arc,
            String, bool, u32, u64, f64, usize},
  keywordstyle=\color{accentblue}\bfseries,
  comment=[l]{//},
  commentstyle=\color{codecomment}\itshape,
  stringstyle=\color{codestring},
  string=[b]",
  sensitive=true,
  morecomment=[s]{/*}{*/},
}

\lstset{
  language=Rust,
  backgroundcolor=\color{codegray},
  basicstyle=\ttfamily\footnotesize,
  breaklines=true,
  keepspaces=true,
  showstringspaces=false,
  frame=single,
  framerule=0.5pt,
  rulecolor=\color{bordercolor},
  xleftmargin=12pt,
  xrightmargin=12pt,
  aboveskip=8pt,
  belowskip=8pt,
  numbers=left,
  numberstyle=\tiny\color{codenumber},
  numbersep=8pt,
  escapeinside={(*@}{@*)},
}

\newcommand{\entity}[1]{\texttt{\textbf{#1}}}
\newcommand{\field}[1]{\texttt{#1}}

\newenvironment{decisionbox}{%
  \begin{mdframed}[backgroundcolor=notegreen,
    linecolor=noteborder, linewidth=1pt,
    innertopmargin=8pt, innerbottommargin=8pt,
    innerleftmargin=10pt, innerrightmargin=10pt]
  \small
}{%
  \end{mdframed}
}

\newenvironment{cautionbox}{%
  \begin{mdframed}[backgroundcolor=warningyellow,
    linecolor=warningborder, linewidth=1pt,
    innertopmargin=8pt, innerbottommargin=8pt,
    innerleftmargin=10pt, innerrightmargin=10pt]
  \small
}{%
  \end{mdframed}
}

\title{Prescription Modeling in the FEN Schema:\\
Facts, ProblemEpisodes, and Medication-State Assessment}
\author{Phoros Engineering}
\date{August 2026}

\begin{document}
\maketitle

\section{Status and Adopted Direction}

This document specifies how prescription and medication-order information
participates in the Facts--Episodes--Narratives (FEN) ontology. It replaces
the earlier comparison among four mutually exclusive lifecycle options.

\begin{decisionbox}
\textbf{Adopted direction.}

\begin{enumerate}[noitemsep]
  \item A medication order and each subsequent lifecycle event are Facts.
  \item Medication Facts may be grouped through authored
    \entity{EpisodeMembership} records into a medication-management
    \entity{ProblemEpisode}.
  \item The schema does not determine the boundary of that
    \entity{ProblemEpisode}. The boundary is established by whoever authors
    the memberships and remains contestable in the ordinary FEN manner.
  \item Source medication-order identifiers are external references. They
    may support an author's membership judgment, but they are not FEN course
    identifiers and do not determine episode identity.
  \item Current medication state is assessed from applicable Facts. A
    transient assessment is a projection. If Phoros persists the assessment
    and takes epistemic responsibility for it, the assessment is a
    Phoros-authored inference Fact whose assertion method is
    \field{Inferred} and whose evidence cites its inputs.
\end{enumerate}
\end{decisionbox}

No fourth top-level medication entity is introduced. A mutable
\entity{MedicationOrder} may still be useful as an interoperability resource,
API representation, or materialized read model, but it is not part of the
core FEN ontology.

\subsection{Reading the source-to-Fact boundary}

Suppose an EHR export contains a medication order in a PDF or EHI table. The
received file is preserved outside FEN as an immutable
\entity{SourceArtifact}, with its bytes held by a \entity{StoredObject}. A
\entity{SourceCitationRef} identifies the relevant page, region, table row, or
field. The normalized medication event is a Fact: its
\field{assertion: AssertionProvenance} records who or what admitted the claim,
when, and how; its \field{evidence: Vec<EvidenceRef>} cites the source location
or prior Facts that support it. A source order identifier belongs in
\field{external\_refs}; it supports correlation, not epistemic justification.

\begin{center}
\small
\renewcommand{\arraystretch}{1.12}
\begin{tabular}{>{\raggedright\arraybackslash}p{4.8cm} >{\raggedright\arraybackslash}p{8.0cm}}
\toprule
\textbf{Concern} & \textbf{Representation} \\
\midrule
Original medication document or export & \entity{SourceArtifact} $\rightarrow$ \entity{StoredObject} \\
Exact supporting passage or row & \entity{SourceCitationRef} \\
Who asserted the normalized event and how & \entity{AssertionProvenance} \\
Supporting source material or prior Facts & \field{Fact.evidence} \\
Source order, request, version, or event identity & \field{Fact.external\_refs} \\
Transient processing or display state & \entity{FactOperational} or projection \\
\bottomrule
\end{tabular}
\end{center}

\section{Epistemic Commitments}
\label{sec:epistemic-commitments}

\subsection{Facts carry assertion metadata and evidence}

A Fact is not guaranteed truth. It is an epistemic record admitted to the
system with assertion metadata, evidence, a temporal anchor, and status.
Assertion metadata records who or what asserted the Fact, when, and by which
method. Evidence records the source citations and prior Facts that support the
claim. Medication
Facts may be directly observed, authored by a clinician or patient, reported
by an institution, extracted from a source document, or inferred by Phoros.
These distinctions determine how a Fact should be evaluated; they do not
determine whether it is ontologically a Fact.

Examples include:

\begin{itemize}[noitemsep]
  \item a clinician authored a medication order;
  \item an EHR reported that an order was modified or discontinued;
  \item a pharmacy reported that medication was dispensed;
  \item a patient stated that they stopped taking a medication;
  \item Phoros inferred a current medication state or a possible interaction;
  \item a clinician authored a Note explaining the reason for a change.
\end{itemize}

A clinician-authored Note is therefore prime Fact material. Its content may be
challenged, contradicted, refined, or interpreted by other epistemic objects,
but the event and content of its authorship remain a Fact.

\subsection{ProblemEpisodes are authored problem-oriented groupings}

A \entity{ProblemEpisode} is used when treating a collection of Facts
separately obscures a problem those Facts address together. A medication
course can be represented as such a problem-oriented grouping, but the schema
does not supply an objective course-boundary algorithm.

An author may judge that two renewals belong to one medication-management
ProblemEpisode, while another author may distinguish them. One source order
may participate in multiple ProblemEpisodes. Multiple source order identifiers
may participate in one ProblemEpisode. A medication Fact may initially have
no Episode membership at all.

All such cases are represented by ordinary \entity{EpisodeMembership}
assertion metadata, status, and contestability. Membership is not a Fact, so
its \field{asserted\_by} and \field{asserted\_at} fields remain distinct from
\field{Fact.assertion}. A source order identifier, shared ingredient code,
temporal adjacency, or automated correlation result may support a membership
judgment; none is a schema-level rule that creates the boundary. If membership
evidence becomes first-class, it must remain distinct from the metadata about
who authored the membership.

\subsection{Narratives provide higher-order accounts}

A \entity{Narrative} is a higher-order authored account of a ProblemEpisode.
It is not distinguished from Fact merely because it contains interpretation:
authored and inferred assertions may also be Facts. Narrative is distinguished
by its episode-level explanatory and synthesizing role.

For example, a clinician's Note explaining a dose change is a Fact. A later
longitudinal account explaining how renal function, adverse effects, and
patient preferences shaped several medication decisions is a Narrative.

\subsection{Projections are not a fourth ontological layer}

Read models, indexes, and transient calculations may be materialized for
performance or presentation without becoming new FEN entity types. The
distinction is one of epistemic commitment:

\begin{itemize}[noitemsep]
  \item A transient calculation used to render a screen is a projection.
  \item A persisted assertion for which Phoros records authorship and
    evidence is an inference Fact with \field{AssertionMethod::Inferred}.
\end{itemize}

\section{Problems with the Prior Prescription Payload}
\label{sec:current-problems}

The prior payload was approximately:

\begin{lstlisting}
Prescription {
    medication:    CodedValue,
    dosage:        Dosage,
    frequency:     Frequency,     // previously Vec<Duration>
    period:        Option<TimeInterval>,
    relevant_note: Option<FactId>,
}
\end{lstlisting}

It has five separable problems.

\subsection{Frequency cannot be represented as \texttt{Vec<Duration>}}

A sequence of time gaps cannot faithfully express common medication
instructions. It cannot adequately distinguish ``twice daily'' from ``every
twelve hours,'' or represent PRN use, frequency ranges, day-of-week schedules,
times of day, relationships to meals, titration stages, conditional holds,
rates, or maximum doses.

Medication instructions must preserve both structured information and the
original SIG text. The exact Rust representation can evolve independently of
the lifecycle decision, but it must support at least:

\begin{itemize}[noitemsep]
  \item ordered or concurrent instruction stages;
  \item frequency and frequency ranges over a period;
  \item days of week, times of day, and event-relative timing;
  \item PRN indication and other conditions;
  \item dose, route, method, rate, and maximum-dose constraints;
  \item intended bounds; and
  \item original source text and source coding.
\end{itemize}

A representative internal shape is:

\begin{lstlisting}
pub struct MedicationInstruction {
    pub dosage: Dosage,
    pub timing: MedicationTiming,
    pub route: Option<CodedValue>,
    pub method: Option<CodedValue>,
    pub as_needed_for: Vec<CodedValue>,
    pub maximum_dose: Option<MaximumDose>,
    pub original_sig: Option<String>,
}

pub struct MedicationTiming {
    pub code: Option<CodedValue>,
    pub repeat: Option<RepeatPattern>,
    pub event_offsets: Vec<EventOffset>,
    pub text: Option<String>,
}
\end{lstlisting}

This is a requirements sketch rather than a final medication-timing
specification.

\subsection{Temporal fields had overlapping meanings}

The former \field{Fact.occurred\_at} and
\field{Prescription.period} fields did not state whether they meant order
authorship, intended effectiveness, actual medication use, or the duration of
a retrospectively grouped course.

The revised contract is:

\begin{itemize}[noitemsep]
  \item \field{Fact.occurred\_at} records when the epistemic event represented
    by the Fact occurred or was reported to occur.
  \item \field{effective\_from} records a source-authored prospective
    effective date when it differs from the order event.
  \item \field{intended\_duration} records the stated prospective intent.
  \item A stop, hold, resume, or modification is represented by another Fact.
  \item A ProblemEpisode does not acquire a schema-determined interval merely
    because its medication Facts can be temporally ordered.
\end{itemize}

An inferred actual duration may be calculated from selected Facts and
memberships, but it remains an assessment whose result can be ambiguous.

\subsection{The lifecycle event was implicit}

The former payload could show two different medication snapshots but could
not state whether the difference represented an initial order, correction,
dose modification, hold, renewal, or discontinuation. Comparing values is not
a substitute for preserving the event asserted by the source.

\subsection{Fact status was easily confused with medication state}

A later dose change does not make the earlier order cease to have occurred.
Likewise, discontinuation does not make prior order and modification Facts
epistemically inactive. Medication lifecycle state and \entity{FactStatus}
answer different questions and must not share one state machine.

\subsection{One relevant Note reference was too narrow}

A medication assertion may be supported by Notes, laboratory results,
diagnoses, dispensing records, or Facts from another domain. The singular
\field{relevant\_note} field is therefore removed. General epistemic support
belongs in the enclosing Fact's evidence collection. Prior Facts use the
\field{PriorFact} evidence variant; source passages use \field{Source}. A
referenced Fact retains its own authorization policy; the reference does not
grant access to its target.

\section{Adopted Medication Event Fact Model}
\label{sec:adopted-model}

The adopted design combines the best properties of the earlier event-stream
and order-plus-delta proposals. Every lifecycle event is a Fact. An initial
order carries a complete instruction. A modification records what the source
says changed and, when the source supplies it, the complete resulting
instruction.

\subsection{Payload family}

Medication payloads belong to the closed clinical Fact payload family. The
shared FEN ontology does not require one global enum containing identity,
clinical, and health-economic variants.

\begin{lstlisting}
pub enum ClinicalFactPayload {
    // Other closed clinical variants...
    MedicationOrderEvent(MedicationOrderEventPayload),
    MedicationCourseAssessment(MedicationCourseAssessmentPayload),
    // The interaction representation remains deferred.
}

pub struct MedicationOrderEventPayload {
    pub medication: CodedValue,
    pub event: MedicationOrderEvent,
}
\end{lstlisting}

The enclosing Fact supplies the common identity, subject, temporal, status,
assertion, evidence, and external-reference fields. Payload-level Fact IDs are
retained only where they express medication-specific domain semantics, such as
the event targeted by a modification.

\subsection{Lifecycle events}

\begin{lstlisting}
pub enum MedicationOrderEvent {
    Authored {
        instruction: MedicationInstruction,
        intended_duration: Option<Duration>,
        effective_from: Option<ApproximateDate>,
    },

    Modified {
        target_fact_id: Option<FactId>,
        changes: MedicationInstructionPatch,
        resulting_instruction: Option<MedicationInstruction>,
        reason: Option<String>,
    },

    Held {
        target_fact_id: Option<FactId>,
        reason: Option<String>,
        expected_resume: Option<ApproximateDate>,
    },

    Resumed {
        target_fact_id: Option<FactId>,
    },

    Discontinued {
        target_fact_id: Option<FactId>,
        reason: Option<String>,
    },

    Renewed {
        prior_order_fact_id: Option<FactId>,
        instruction: MedicationInstruction,
        intended_duration: Option<Duration>,
    },

    ReportedSnapshot {
        instruction: MedicationInstruction,
        reported_status: Option<ReportedMedicationStatus>,
        effective_from: Option<ApproximateDate>,
        effective_through: Option<ApproximateDate>,
    },
}
\end{lstlisting}

\field{ReportedSnapshot} is retained for sources and migrations that provide
a medication state without reliable evidence of the lifecycle event that
produced it. Ingestion must not invent \field{Authored} or \field{Modified}
semantics merely to fit a preferred event model.

\subsection{Patch semantics}

A modification must distinguish an omitted field from a field explicitly
removed by the source.

\begin{lstlisting}
pub struct MedicationInstructionPatch {
    pub dosage: Option<FieldChange<Dosage>>,
    pub timing: Option<FieldChange<MedicationTiming>>,
    pub route: Option<FieldChange<CodedValue>>,
    pub method: Option<FieldChange<CodedValue>>,
    pub as_needed_for: Option<FieldChange<Vec<CodedValue>>>,
    pub maximum_dose: Option<FieldChange<MaximumDose>>,
    pub original_sig: Option<FieldChange<String>>,
}

pub enum FieldChange<T> {
    Set(T),
    Remove,
}
\end{lstlisting}

The two fields on \field{Modified} make different epistemic commitments:

\begin{itemize}[noitemsep]
  \item \field{changes} records what the source identified as changing.
  \item \field{resulting\_instruction} records a complete post-change
    instruction only when the source explicitly supplies one.
\end{itemize}

Phoros must not interpret an omitted field as a source-authored assertion that
the field remained unchanged unless the source contract establishes that
semantics.

\subsection{Source identifiers}

The model introduces no FEN \field{MedicationCourseId} or internal
\field{OrderGroupId}. Source order, request, version, and event identifiers
are recorded through \field{Fact.external\_refs}. The
\field{resource\_type} or equivalent discriminator distinguishes an order
identifier from an event or version identifier.

These identifiers answer ``which source record is this?'' They do not answer
``what supports this assertion?'' The latter requires a
\entity{SourceCitationRef} to the relevant source passage, row, or field in
\field{Fact.evidence}.

A \field{target\_fact\_id} is useful when a source or author states that one
event specifically changes another. It is not a substitute for
ProblemEpisode membership and is not required when the target cannot be
established.

\section{ProblemEpisode Membership}
\label{sec:problem-episode-membership}

Medication Facts participate in the existing FEN membership mechanism:

\begin{lstlisting}
EpisodeMembership {
    id: MembershipId("mem-2"),
    fact_id: FactId("f-metformin-modified"),
    episode_id: ProblemEpisodeId("ep-metformin-management"),
    role: FactRole::Treatment,
    asserted_by: Author::clinician("Dr. Patel"),
    asserted_at: TemporalAnchor::Point(Timestamp("2026-01-16")),
    status: MembershipStatus::Active,
}
\end{lstlisting}

The schema deliberately does not answer whether a renewal, restart, changed
indication, formulation change, or long treatment gap begins a new
ProblemEpisode. That is a question for the membership author. Automated
systems may author memberships, but they must do so with identifiable
authorship and the same contestability expected of other membership
assertions. Automated authoring also records the rule or model version in the
appropriate audit or operational record; this mechanism metadata must not be
collapsed into membership authorship.

A medication Fact may also belong to a broader clinical ProblemEpisode. For
example, the same dose-modification Fact may be a member of both a metformin
management ProblemEpisode and a chronic kidney disease ProblemEpisode, with a
different \entity{FactRole} in each.

\section{Fact Status and Medication Lifecycle State}
\label{sec:status-distinction}

\entity{FactStatus} expresses the system's current disposition toward an
epistemic record:

\begin{lstlisting}
pub enum FactStatus {
    Active,
    Superseded {
        replaced_by: Option<FactId>,
        reason: SupersessionReason,
        // author and time omitted here for brevity
    },
    EnteredInError {
        replaced_by: Option<FactId>,
        // author and time omitted here for brevity
    },
}
\end{lstlisting}

Medication lifecycle state expresses what the available evidence says about
the medication:

\begin{lstlisting}
pub enum MedicationLifecycleState {
    Ordered,
    Active,
    OnHold,
    Discontinued,
    Completed,
    Unknown,
}
\end{lstlisting}

These states must not be conflated. Consider:

\begin{center}
\renewcommand{\arraystretch}{1.25}
\begin{tabular}{p{6.4cm} p{3.2cm} p{3.2cm}}
\toprule
\textbf{Fact} & \textbf{Fact status} & \textbf{Course effect} \\
\midrule
Metformin 500 mg ordered on January 1 & Active & Establishes order \\
Dose increased to 1000 mg on March 1 & Active & Changes dose \\
Medication discontinued on June 1 & Active & Ends active use \\
\bottomrule
\end{tabular}
\end{center}

All three Facts remain valid records of events that occurred even though the
assessed medication state is \field{Discontinued}.

Supersession is appropriate when a new Fact replaces the epistemic position
of an earlier Fact. For example, an imported Fact saying that the original
order was 500 mg may be superseded by a corrected Fact grounded in
authoritative source evidence that establishes it was 250 mg. Ordinary dose
changes, holds, renewals, and
discontinuations are not supersessions.

\field{Fact.status} may store the materialized current disposition. The
history of status transitions must nevertheless remain append-only and
reconstructable so the system can answer both ``what is the Fact's status
now?'' and ``what did the system regard as active at time $t$?'' The storage
mechanism for that history is outside this prescription specification.

\section{Medication-Course State Assessment}
\label{sec:state-assessment}

The earlier term \field{FoldMedicationState} implied a complete, consistent,
totally ordered event stream. FEN cannot assume those conditions. Relevant
Facts may conflict, come from different sources, have uncertain membership,
or support more than one present-state interpretation.

The evaluator is therefore described as an assessment:

\begin{lstlisting}
pub enum MedicationCourseAssessment {
    Determinate {
        state: CurrentMedicationState,
    },

    Ambiguous {
        candidate_states: Vec<CurrentMedicationState>,
        conflicting_facts: Vec<FactId>,
    },

    InsufficientEvidence {
        relevant_facts: Vec<FactId>,
    },
}
\end{lstlisting}

An assessment request identifies a ProblemEpisode, an \field{as\_of} time,
and an authorized view of its memberships and Facts. The evaluator must
record the rule or model version and all input Fact IDs on which its result
depends.

If the result is transient, it remains a projection. If Phoros persists the
result as an epistemic assertion, it is represented as a Phoros-authored
inference Fact:

\begin{lstlisting}
pub struct MedicationCourseAssessmentPayload {
    pub problem_episode_id: ProblemEpisodeId,
    pub as_of: Timestamp,
    pub assessment: MedicationCourseAssessment,
    pub evaluator: RuleArtifactRef,
}
\end{lstlisting}

The enclosing inference Fact records \field{AssertionMethod::Inferred} and
cites every epistemic input through \field{Fact.evidence}. The
\field{conflicting\_facts} and \field{relevant\_facts} fields remain inside the
assessment because they explain why the result is ambiguous or insufficient;
when the assessment is persisted, those same Fact IDs must also appear as
evidence. This deliberate duplication is validated for consistency.

Persisted assessments may be superseded when new evidence arrives or the
evaluator changes. That is epistemic replacement and is appropriately
represented through \entity{FactStatus} and
\field{SupersessionReason::RuleReEvaluation}.

\section{Notes, Evidence, and Inferred Rationales}
\label{sec:notes-and-rationales}

An event payload may carry a reason when that reason was part of the source
assertion. A clinician-authored Note that explains the decision is itself a
Fact and may be cited through \field{Fact.evidence}. The medication event's
source citation establishes what the clinician actually authored.

Phoros must not insert an inferred rationale into a clinician-authored Fact
as though the clinician supplied it. Instead, Phoros authors a separate
inference Fact:

\begin{lstlisting}
pub struct MedicationRationaleInferencePayload {
    pub medication_event_fact_id: FactId,
    pub rationale: String,
    pub confidence: Option<Confidence>,
    pub evaluator: RuleArtifactRef,
}
\end{lstlisting}

This permits a Phoros inference to cite the medication event, relevant Note,
laboratory results, and Facts from other domains without changing the
authorship of any original Fact. The payload's
\field{medication\_event\_fact\_id} identifies the assertion being explained;
the enclosing inference Fact's \field{evidence} records the inputs that support
the inferred rationale.

\section{Medication Interactions}
\label{sec:interactions}

Medication ``interaction'' covers several distinct claims:

\begin{enumerate}[noitemsep]
  \item contraindication or conflict;
  \item counter-medication used to manage an effect of another medication;
  \item intentional joint regimen;
  \item dose-modifying interaction;
  \item causal chain in which one medication creates a need for another; and
  \item sequential replacement or supersession for a stated reason.
\end{enumerate}

General drug--drug knowledge is contextual reference data, not a patient Fact.
When a clinician, institution, knowledge system, or Phoros asserts that an
interaction applies to a particular subject, that patient-specific assertion
may be represented as a Fact:

\begin{lstlisting}
pub struct MedicationInteractionAssertionPayload {
    pub involved_facts: Vec<FactId>,
    pub interaction_type: MedicationInteractionType,
    pub assertion: Option<String>,
    pub knowledge_source: Option<ReferenceDataRef>,
    pub evaluator: Option<RuleArtifactRef>,
    pub confidence: Option<Confidence>,
}

pub enum MedicationInteractionType {
    Contraindication,
    CounterMedication,
    JointRegimen,
    DoseModifying,
    CausedNeedFor,
    Replaced,
}
\end{lstlisting}

The enclosing Fact records assertion metadata and evidence. The
\field{involved\_facts} field identifies the medications participating in the
interaction; it is domain content rather than a substitute for evidence. An
inferred interaction cites its patient-specific input Facts and the relevant
location in the versioned knowledge artifact. \field{ReferenceDataRef}
identifies the knowledge version, while a source citation identifies the
supporting passage; \field{evaluator} identifies the rule or model that applied
it.

Other compatible representations remain available for different purposes:

\begin{itemize}[noitemsep]
  \item A typed relation between ProblemEpisodes may express a relationship
    between two medication-management problems.
  \item \entity{EpisodeMembership.role} may state the function a Fact plays
    inside one ProblemEpisode, but cannot by itself represent a complete
    pairwise interaction.
  \item A Narrative or DecisionPoint may explain how an interaction affected
    treatment reasoning.
\end{itemize}

This document does not yet choose one authoritative representation for every
interaction use case. The adopted medication lifecycle model is compatible
with all three and does not require a fourth medication entity.

\section{Rejected Alternatives}
\label{sec:rejected-alternatives}

\subsection{Using Fact supersession for clinical evolution}

Rejected. A dose modification, hold, resume, renewal, or discontinuation is a
new epistemic event, not a correction of the earlier event. Reusing
\entity{FactStatus::Superseded} would erase the distinction between change in
the world and replacement of an epistemic record.

\subsection{A mandatory internal medication-order grouping identifier}

Rejected as the FEN course identity. Source order identifiers remain external
references, and ProblemEpisode membership supplies authored grouping. An
operational correlation key may exist in an ingestion or read-model layer,
but it does not determine ProblemEpisode boundaries.

\subsection{A separate top-level MedicationOrder entity}

Rejected as a core FEN entity. The proposed mutable object combined event
history, current state, and grouping into a fourth semantic layer. Its useful
features are already covered by medication event Facts, ProblemEpisodes, and
state assessments. A derived MedicationOrder representation remains
permissible for FHIR interchange, APIs, or optimized reads.

\subsection{Full snapshots only}

Rejected as the sole lifecycle representation. Full snapshots are retained
when that is all a source supplies, but they must be labeled as reported
snapshots. When the source explicitly reports a modification, the Fact should
preserve the modification semantics rather than requiring consumers to infer
them by comparing snapshots.

\section{Scope Beyond Prescription}

Coverage, prior authorization, diagnosis certainty, and other longitudinal
domains may exhibit a similar pattern: multiple Facts become intelligible
only when considered under a ProblemEpisode. Their specific payloads,
membership practices, and state assessments belong to their respective
domain specifications. This prescription document makes no binding lifecycle
decision for them.

\section{Migration Considerations}
\label{sec:migration}

Existing prescription records must be migrated without inventing epistemic
specificity that the source does not contain.

Migration begins by registering the original document, archive, or export as a
\entity{SourceArtifact} and preserving its bytes in a \entity{StoredObject}.
The migration then creates citations to the relevant medication-list row,
order, or Note before constructing normalized Facts.

\begin{enumerate}
  \item An existing Fact with evidence of an initial authored order maps to
    the \field{Authored} event variant.
  \item An existing source snapshot whose lifecycle event cannot be
    established maps to \field{ReportedSnapshot}.
  \item Existing source order identifiers migrate to
    \field{Fact.external\_refs}; they do not become ProblemEpisode IDs.
  \item The old relevant-Note value becomes a \field{PriorFact} entry in
    \field{Fact.evidence}.
  \item Existing Facts are not assigned to ProblemEpisodes by schema
    migration alone. A separate authored or automated membership process may
    propose memberships with explicit authorship and audit metadata.
  \item Existing Facts are not superseded merely because a later medication
    record exists. Supersession is preserved only where the migration can
    establish genuine epistemic replacement or correction.
  \item A migration-generated resulting instruction is an inferred assertion.
    It cites the source material and prior Facts used to build it and records
    the reconstruction rule version.
\end{enumerate}

This ordering ensures that normalization never becomes the sole surviving
representation of the source and that future parsers can revisit the original
meaning.

\section{Decision Summary}

\begin{center}
\footnotesize
\renewcommand{\arraystretch}{1.08}
\begin{tabular}{>{\raggedright\arraybackslash}p{5.0cm} >{\raggedright\arraybackslash}p{8.0cm}}
\toprule
\textbf{Concern} & \textbf{FEN representation} \\
\midrule
Initial medication order & Fact \\
Modification, hold, resume, stop, renewal & Facts \\
Original medication document or export & \entity{SourceArtifact} outside FEN \\
Exact supporting passage, row, or field & \entity{SourceCitationRef} \\
Supporting Note, lab, or prior medication Fact & \field{EvidenceRef::PriorFact} \\
Source medication-order identity & \field{ExternalRef} \\
Medication-management grouping & \entity{ProblemEpisode} plus authored \entity{EpisodeMembership} \\
ProblemEpisode boundary & Determined by membership authors, not by schema \\
Specific event-to-event dependency & Direct \field{FactId} or typed Fact relation \\
Present-state calculation & Projection or Phoros-authored inference Fact \\
Clinician Note and stated rationale & Fact \\
Phoros-inferred rationale & Separate inference Fact \\
Correction of an epistemic record & \entity{FactStatus::Superseded} or \entity{EnteredInError} \\
Medication no longer active & Medication lifecycle assessment, not \entity{FactStatus} \\
Mutable interoperability/API resource & Derived or materialized representation outside core FEN \\
Transient processing or indexing state & \entity{FactOperational} or projection \\
\bottomrule
\end{tabular}
\end{center}

{\small The central commitment is that prescription modeling uses the existing FEN
ontology without collapsing clinical evolution into Fact replacement and
without introducing a fourth medication layer. Facts preserve what was
authored, observed, reported, or inferred. ProblemEpisodes group those Facts
under author-defined problems. Narratives synthesize their meaning. Read
models and transient assessments remain operational artifacts unless Phoros
persists them as inference Facts with explicit assertion metadata and evidence.}

\end{document}
